
\section{Electricity selling price}
\label{app:selling_price}

The price received for exported electricity applies the spot price less a marketing spread,
floored, net of a self-consumption share and a fixed fee, and cannot go negative:
\begin{equation}
  \lambda_t^{\mathrm{sell}} = \max\!\Bigl(
    \max\bigl(\lambda_t^{\mathrm{spot}} - \lambda^{\mathrm{spread}},\;
      \lambda^{\mathrm{floor}}\bigr)
    (1 - h^{\mathrm{sell}})
    - \lambda^{\mathrm{sell,fee}},\; 0\Bigr).
\end{equation}

\section{Per-node temperature error breakdown}
\label{app:per_node_mae}

Table~\ref{tab:per_node_mae} reports supply-temperature error at all fourteen network nodes
over a 744-hour winter window, evaluated on the temperature-propagating formulation and
sorted by path length from the source. Three groups are distinguished, and the distinction
matters for how the aggregate figures in Section~\ref{subsec:validation} should be read.

\emph{Validation nodes} (six) are consumer sensors whose mixing-valve offset is below
$2$\,°C and which took no part in the calibration; their mean error is $1.32$\,°C.
\emph{Calibration nodes} (four) were used in the branch-sequential coefficient tuning, so
their residuals are in-sample by construction; their mean is $2.32$\,°C, and one node at
$3.07$\,°C would fail a $2.5$\,°C held-out criterion. That in-sample deterioration is
expected, and it is the reason a strict calibration/validation split matters.
\emph{Excluded nodes} (four) are set aside for measurement reasons unrelated to model
quality: two exhibit mixing-valve absorption above the $2$\,°C threshold and cannot be read
as junction temperatures, one is a boundary-condition reference, and one serves four
consumers with a $20.9$\,°C intra-node spread, so no single junction reference exists.

The mean across all fourteen nodes is $1.68$\,°C, between the validation-node mean of
$1.32$\,°C and the calibration-node mean of $2.32$\,°C. The validation nodes are therefore
not systematically the best-performing ones (two of the excluded nodes are in fact the
best in the network), so the group means are not an artefact of favourable node selection.

\medskip
{%
\captionof{table}{Per-node supply-temperature error over a 744\,h winter window, evaluated
on the temperature-propagating formulation.}
\label{tab:per_node_mae}
\centering
\footnotesize
\begin{threeparttable}
\begin{tabular*}{\columnwidth}{@{\extracolsep{\fill}} l c c c c c @{}}
\toprule
Node & Path [m] & Role & MAE [$^\circ$C] & Bias [$^\circ$C] & In gate? \\
\midrule
j$_2$    &  350 & excl. & 2.99 & $+2.98$ & No  \\
j$_3$    &  650 & excl. & 0.58 & $+0.07$ & No  \\
j$_4$    &  910 & excl. & 0.37 & $-0.18$ & No  \\
j$_9$    & 1100 & val.  & 0.87 & $+0.41$ & Yes \\
j$_5$    & 1150 & excl. & 2.35 & $+1.51$ & No  \\
j$_6$    & 1300 & cal.  & 1.62 & $-1.62$ & No  \\
j$_{10}$ & 1300 & val.  & 2.27 & $+1.78$ & Yes \\
j$_7$    & 1370 & cal.  & 3.07\tnote{a} & $+2.53$ & No  \\
j$_8$    & 1450 & cal.  & 2.29 & $+1.00$ & No  \\
j$_{11}$ & 1530 & val.  & 1.41 & $+0.71$ & Yes \\
j$_{12}$ & 1780 & val.  & 0.69 & $-0.29$ & Yes \\
j$_{13}$ & 2000 & val.  & 1.09 & $-1.02$ & Yes \\
j$_{15}$ & 2125 & val.  & 1.58 & $-0.24$ & Yes \\
j$_{14}$ & 2180 & cal.  & 2.31 & $+1.02$ & No  \\
\midrule
\multicolumn{3}{@{}l}{Group mean, validation (6 nodes)}  & 1.32 & & \\
\multicolumn{3}{@{}l}{Group mean, calibration (4 nodes)} & 2.32 & & \\
\multicolumn{3}{@{}l}{Group mean, excluded (4 nodes)}    & 1.57 & & \\
\multicolumn{3}{@{}l}{All 14 nodes}                       & 1.68 & & \\
\bottomrule
\end{tabular*}
\begin{tablenotes}\scriptsize
  \item[a] In-sample residual on a calibration node; would exceed a $2.5$\,°C held-out
    criterion.
\end{tablenotes}
\end{threeparttable}
}% end centering group

\section{Demand heterogeneity index}
\label{app:hi_definition}

With $d_n$ the share of total annual demand at node $n$, the normalised variance index is
\begin{equation}
  d_n = \frac{\sum_t Q_{n,t}^{\mathrm{dem}}}{\sum_n \sum_t Q_{n,t}^{\mathrm{dem}}},
  \qquad
  \mathrm{HI} = \frac{N}{N-1} \sum_{n=1}^{N}
    \left(d_n - \frac{1}{N}\right)^{2}.
  \label{eq:hi}
\end{equation}
The index maps monotonically to the Gini coefficient on tree topologies but compresses the
realistic range: the industrial case has $\mathrm{HI} = 0.05$ against a Gini of $0.41$. The
synthetic values $0.1$, $0.4$ and $0.8$ correspond approximately to Gini ranges of
$0.15$--$0.25$, $0.40$--$0.55$ and $0.70$--$0.85$. Because of that compression, the
model-selection observations in Table~\ref{tab:criteria} are expressed in terms of the Gini
coefficient and the top-$k$ demand share rather than this index.

